\section{Physical question and equivalence conventions}
\label{sec:intrinsic-question}

A boundary density operator can contain local spectral weights and a
global constraint. The weights affect entanglement spectra and boundary
entropies. The constraint can retain information that disappears upon
tracing even one boundary site. The boundary discussion following the
general fixed-point characterization of Cirac, P\'erez-Garc\'ia,
Schuch, and Verstraete (CPSV)~\cite{CPSV} motivates asking whether
length-dependent matrix product density operator (MPDO) fixed points
can involve more than weighted dimers.

There are three different questions. Can canonical structure
coefficients depend on length at a specified physical cut? Can such
an example have nontrivial fusion and a state that is not an independent
dimer product? Can its positive weights remain inseparable from its
intertwiners under a specified class of local operations? The examples
below answer the first two questions affirmatively. They do not yet
settle the third under bounded-depth local refactoring.

This distinction matters physically. Stacking a boundary with a trivial
layer of mixed bonds changes its entanglement spectrum without changing
the topological order of an associated bulk. Conversely, identifying
weights in local multiplicity spaces need not produce an independent
tensor factor. The categorical examples make the latter point explicit.

\subsection{Tensor orientations and physical normalization}

Let $\mc H$ be a finite-dimensional physical site space and
$\mc V=\C^D$ the horizontal auxiliary space. A tensor is a collection
$M^{ab}\in\End(\mc H)$, with $a,b=1,\ldots,D$. Define
\begin{align}
	\rho_N(M)
	 & =\sum_{a_1,\ldots,a_N}
	M^{a_1a_2}\otimes\cdots\otimes M^{a_Na_1},
	\label{eq:horizontal}            \\
	(MB)^{ab}
	 & =\sum_c M^{ac}\otimes B^{cb},
	\label{eq:concatenation}         \\
	[O_L(M)]_{\boldsymbol a,\boldsymbol b}
	 & =\tr_{\mc H}
	\left(M^{a_1b_1}\cdots M^{a_Lb_L}\right).
	\label{eq:vertical}
\end{align}
Here $N$ counts physical sites, whereas $L$ counts factors in the
vertical trace. The two readings use the same tensor coefficients,
without transposition or complex conjugation. Positivity of every
$\rho_N(M)$ does not imply positivity of $O_L(M)$.

A physical renormalization fixed point (RFP) has completely positive,
trace-preserving maps, abbreviated CPTP maps,
\begin{align}
	\mc T:\End(\mc H)            & \longrightarrow\End(\mc H^{\otimes2}), &
	\mc S:\End(\mc H^{\otimes2}) & \longrightarrow\End(\mc H),
	\nonumber                                                               \\
	\mc T(M^{ab})                & =(MM)^{ab},                            &
	\mc S((MM)^{ab})             & =M^{ab}.
	\label{eq:rfp}
\end{align}
These are identities for every complex letter, not just for positive
periodic states. Where canonical form is required, it is established
separately.

Trace preservation fixes the physical scale. With
$T_{ab}=\tr(M^{ab})$, one obtains
\begin{align}
	T_{ab}
	 & =\tr\mc T(M^{ab})
	=\sum_c\tr(M^{ac})\tr(M^{cb})
	=(T^2)_{ab},                \\
	\tr\rho_N(M)
	 & =\tr(T^N)=\tr T=\rank T.
	\label{eq:physical-scale}
\end{align}
For a nonzero positive family, both $M$ and $sM$, with $s>0$, can
satisfy this necessary condition only if
$(s^2-s)T=0$, hence $s=1$.

For a tuple $A=(A^u)_u$, define its transfer map by
\begin{align}
	\mc E_A(X) & =\sum_uA^uX(A^u)^\dagger.
\end{align}
An irreducible, aperiodic tuple with transfer spectral radius $r>0$
has normal representative $A/\sqrt r$. A basis of normal tensors
(BNT) contains one representative of each similarity-and-phase class
in canonical form. Its periodic vectors are independent at all
sufficiently large lengths. We use unit-transfer normal representatives
but retain the external coefficients needed at the physical scale.
Normalizing a representative is not a rescaling of the physical
tensor.

\subsection{Fusion and the three notions of nontriviality}

A vertical canonical decomposition and its fusion identities have the
form
\begin{align}
	M^{ab}
	 & =V_1\left(\bigoplus_c\mu_c\otimes A_c^{ab}\right)V_1^\dagger,
	\label{eq:vcf}                                                   \\
	(A_aA_b)^{ij}
	 & =W_{ab}\left(\bigoplus_c\chi_{abc}\otimes A_c^{ij}\right)
	W_{ab}^\dagger.
	\label{eq:fusion}
\end{align}
The maps $V_1,W_{ab}$ are isometries from the displayed direct sums
into the corresponding physical spaces. Both identities assert
vanishing on the orthogonal complements and the absence of cross
terms. All positive multiplicity operators are restricted to their
active supports. An absent channel has dimension zero.

The product of vertical closures follows from the definitions:
\begin{align}
	[O_L(A)O_L(B)]_{\boldsymbol i,\boldsymbol j}
	 & =\sum_{\boldsymbol k}
	\tr\left(\prod_{\ell=1}^LA^{i_\ell k_\ell}\right)
	\tr\left(\prod_{\ell=1}^LB^{k_\ell j_\ell}\right)
	\nonumber                                          \\
	 & =\tr\left[\prod_{\ell=1}^L
		       \left(\sum_{k_\ell}A^{i_\ell k_\ell}\otimes
		       B^{k_\ell j_\ell}\right)\right] \\
	 & =[O_L(AB)]_{\boldsymbol i,\boldsymbol j}.
\end{align}
For $\chi=\sum_qx_q\ket q\bra q$, orthogonality gives
\begin{align}
	\prod_{\ell=1}^L(\chi\otimes A^{i_\ell j_\ell})
	 & =\sum_qx_q^L\ket q\bra q\otimes
	A^{i_1j_1}\cdots A^{i_Lj_L},        \\
	O_L(\chi\otimes A)
	 & =\left(\sum_qx_q^L\right)O_L(A).
\end{align}
Consequently,
\begin{align}
	O_L(A_a)O_L(A_b)
	           & =\sum_c c_{abc}(L)O_L(A_c), &
	c_{abc}(L) & =\tr(\chi_{abc}^L).
	\label{eq:power-law}
\end{align}
The power-sum law is universal within the stated CPSV
characterization. A different physical mechanism would mean
inseparable weights and intertwiners, not a violation of
\cref{eq:power-law}.

A compensated representative change is
\begin{align}
	A_a'        & =s_aA_a,                       &
	\mu_a'      & =\mu_a/s_a,                    &
	\chi_{abc}' & =\frac{s_as_b}{s_c}\chi_{abc},
	\qquad s_a>0.
	\label{eq:rescaling}
\end{align}
A positive power sum is constant in $L$ only if all its active
eigenvalues are one, since
\begin{align}
	0
	 & =\sum_qx_q-2\sum_qx_q^2+\sum_qx_q^3
	=\sum_qx_q(1-x_q)^2.
	\label{eq:constant-power}
\end{align}
Unequal positive eigenvalues within a channel therefore obstruct
flattening by \cref{eq:rescaling}. For scalar channels
$\chi_{abc}=\theta_{abc}\Id$, flattening instead requires the
simultaneous equations $s_as_b\theta_{abc}=s_c$.

We distinguish similarity gauges, compensated representative
rescalings, fixed regroupings of neighboring subregisters, and
bounded-depth unitary circuits. The latter have bounded gate range
and a rule independent of $N$. A length-dependent change of basis
among the $O_L(A_a)$ is not a fixed tensor transformation.

For the stronger question, a \emph{spectator separation} means an
all-length factorization, after a specified fixed regrouping and
bounded-depth unitary circuit, into a weight-independent core and an
independent product of weighted one-site or bond states. Fixed pure
ancillas may be added or removed. No symmetry preservation is
implicit. The unitaries may depend on the chosen weights but not
on the chain length through a growing range or depth.

This convention excludes arbitrary state-erasing channels. Reversible
CPTP conversion on a specified family is a different possible
equivalence relation; no invariance under every such conversion is
asserted here.

\subsection{A full-space fixed-point certificate}

The following lemma will certify the regrouped tensors without using
broad equivalences involving commuting Gibbs states or saturation of
mutual information.

\begin{lemma}[Equal-mass certificate]
	\label{lem:channels}
	Suppose the active sector set is nonempty and the one-site and blocked
	letters have decompositions
	\begin{align}
		M^{ab}
		 & =\sum_{c,q}\eta_{cq}J_{1cq}A_c^{ab}J_{1cq}^\dagger,  \\
		(MM)^{ab}
		 & =\sum_{c,q}\zeta_{cq}J_{2cq}A_c^{ab}J_{2cq}^\dagger,
	\end{align}
	where $\eta_{cq},\zeta_{cq}>0$, the embeddings have orthogonal
	ranges, both complements are zero on all letters, and
	\begin{align}
		m_c:=\sum_q\eta_{cq}=\sum_q\zeta_{cq}>0.
		\label{eq:equal-masses}
	\end{align}
	Then \cref{eq:rfp} holds with channels on the complete physical spaces.
\end{lemma}

\begin{proof}
	Put $\mc H_1=\mc H$, $\mc H_2=\mc H^{\otimes2}$,
	$P_i=\sum_{c,q}J_{icq}J_{icq}^\dagger$, and
	$\mc A=\bigoplus_c\End(\mc H_c)$. Choose a density operator
	$\tau\in\mc A$. Define
	\begin{align}
		D_i(X)
		 & =\bigoplus_c\sum_qJ_{icq}^\dagger XJ_{icq}
		+\tr((\Id-P_i)X)\tau,                         \\
		R_1(\oplus_cX_c)
		 & =\sum_{c,q}\frac{\eta_{cq}}{m_c}
		J_{1cq}X_cJ_{1cq}^\dagger,                    \\
		R_2(\oplus_cX_c)
		 & =\sum_{c,q}\frac{\zeta_{cq}}{m_c}
		J_{2cq}X_cJ_{2cq}^\dagger.
		\label{eq:discard-prepare}
	\end{align}
	Compression and preparation are completely positive. If
	$\tau=\sum_tp_t\ket t\bra t$ and $\{\ket z\}$ is an
	orthonormal basis of $\ker P_i$, the complement branch has Kraus
	operators $\sqrt{p_t}\ket t\bra z$. Their squared sum is
	$\sum_{t,z}p_t\ket z\bra z=\Id-P_i$. Thus $D_i$ is trace
	preserving on the full input space. For each input sector $c$, the
	squared Kraus sums of $R_1,R_2$ are respectively
	$\sum_q\eta_{cq}\Id/m_c=\Id$ and
	$\sum_q\zeta_{cq}\Id/m_c=\Id$.

	On the letters, the complement branches vanish and
	\begin{align}
		D_1(M^{ab})
		 & =\bigoplus_cm_cA_c^{ab}
		=D_2((MM)^{ab}).
	\end{align}
	Therefore $\mc T=R_2D_1$ and $\mc S=R_1D_2$ satisfy
	\cref{eq:rfp}.
\end{proof}
